\section{Conclusiones}

El presente trabajo desarrolló y evaluó un flujo aplicado para el pronóstico univariado de demanda eléctrica horaria mediante N-BEATS, utilizando datos públicos de Open Power System Data y una comparación experimental contra una referencia ARIMA. La serie de Austria se preparó bajo un contrato uniforme \texttt{timestamp}/\texttt{y}, lo que permitió integrar la preparación de datos, el entrenamiento del modelo, la generación de pronósticos y el cálculo de métricas bajo un procedimiento reproducible.

Los resultados muestran que N-BEATS obtuvo un desempeño favorable frente a ARIMA en tres de los cuatro horizontes evaluados. En el horizonte diario, ambos modelos presentaron errores cercanos y ARIMA obtuvo un RMSE menor, por lo que no puede afirmarse una ventaja clara de N-BEATS en ese caso. En los horizontes semanal, bisemanal y mensual, N-BEATS redujo MAE, RMSE y MAPE respecto a ARIMA, lo que indica que, bajo el protocolo definido, la arquitectura profunda capturó patrones útiles de la serie histórica. Aun así, estos hallazgos deben interpretarse como evidencia acotada a la serie \texttt{AT} y a los tramos de evaluación seleccionados.

La aplicación de escritorio cumplió el flujo funcional propuesto: prepara un workspace externo, trabaja con series bajo un contrato común, permite configurar y entrenar modelos N-BEATS, conserva artefactos de modelo y registra ejecuciones de pronóstico con sus archivos asociados. Esta contribución complementa la evaluación cuantitativa, ya que transforma el experimento en una herramienta operativa capaz de generar y preservar pronósticos de manera organizada. La separación entre interfaz, servicios, modelos y almacenamiento también favorece la mantenibilidad del sistema y permite extender el flujo en trabajos posteriores.

La principal limitación metodológica del estudio es que la evaluación se realizó con un origen de pronóstico por horizonte. Esto permite una comparación controlada entre modelos sobre los mismos timestamps, pero no mide la variabilidad del desempeño en múltiples periodos. Además, el enfoque fue estrictamente univariado, por lo que no incorpora variables exógenas como temperatura, calendario o indicadores operativos. Estas restricciones son coherentes con el alcance inicial del proyecto, pero deben considerarse al interpretar los resultados y al trasladar el sistema a escenarios de uso más amplios.

Como trabajo futuro, se recomienda ampliar la evaluación hacia esquemas multi-origen, incorporar más zonas de Open Power System Data y explorar modelos con variables exógenas cuando existan datos confiables. También sería conveniente fortalecer la trazabilidad del entrenamiento mediante la conservación explícita del subconjunto utilizado o mediante mecanismos de versionado de datos. Estas extensiones permitirían evaluar con mayor robustez la estabilidad del desempeño de N-BEATS y mejorar la reproducibilidad del flujo completo.

En síntesis, el proyecto demuestra que N-BEATS puede integrarse en una aplicación reproducible para pronóstico de demanda eléctrica y ofrecer resultados competitivos frente a una referencia estadística tradicional en los tramos evaluados. La contribución principal no radica únicamente en el uso de una arquitectura profunda, sino en articularla con un flujo de software que organiza datos, modelos, metadatos y pronósticos bajo criterios explícitos de trazabilidad.
